\ssr{ЗАКЛЮЧЕНИЕ}

В процессе выполнения лабораторной работы поставленная цель была успешно достигнута: осуществлен теоретический и практический сравнительный анализ ресурсных затрат и временной сложности алгоритма полного перебора и трех различных реализаций бинарного поиска.

Все запланированные задачи полностью решены:
\begin{enumerate}
	\item описана решаемая согласно варианту задача;
	\item разработаны рекурсивный и нерекурсивные алгоритмы решения задачи, согласно варианту;
	\item описаны разработанные алгоритмы схемами алгоритмов;
	\item описаны требования к программному обеспечению;
	\item описаны средства реализации алгоритмов;
	\item реализованы разработанные алгоритмы;
	\item выполнено тестирование реализации алгоритмов;
	\item выполнена теоретическая оценка затрачиваемой реализацией каждого алгоритма памяти (для рекурсивного алгоритма~--- на материале анализа высоты дерева рекурсивных вызовов и оценки затрачиваемой на один вызов функции памяти);
	\item оценена временная эффективность алгоритмов или их реализаций в лучшем и худшем случае в терминах числа сравнений искомой величины с другими величинами;
	\item представлены результаты сравнительного анализа в формате таблицы;
	\item сделаны выводы из сравнительного анализа реализации разработанных алгоритмов решения одной и той же задачи, заданной вариантом, по критериям временной и ёмкостной эффективности, а также дополнительных требований, налагаемых алгоритмом.
\end{enumerate}

В результате исследования установлено, что выбор метода поиска определяется соотношением между частотой поисковых запросов и затратами на поддержание упорядоченности данных. Для однократных запросов на произвольных коллекциях предпочтителен полный перебор, не требующий накладных расходов на сортировку. Для регулярных операций над структурированными данными неоспоримым преимуществом обладает бинарный поиск.